\chapter*{Preface}
\addcontentsline{toc}{chapter}{Preface}

This work does not begin with a theory to be tested, nor with a question
posed by a textbook. It begins with a table.

Specifically: with the values
\[
1,\quad 11,\quad 20,\quad 46,\quad 63,\quad 105,\quad 130,\quad 188,\quad \ldots
\]
These numbers appeared in a computation. They did not seem to form any
known sequence---at least not one that I was aware of at the time.

What drew my attention was the pattern. The differences between consecutive
values---$10, 9, 26, 17, 42, 25, \ldots$---exhibited a regularity that
could not be ignored. Two subsequences alternate, each arithmetic in its
own right---and their pairwise sums $19, 43, 67, \ldots$ again form an
arithmetic progression. The question of \emph{why} this is so would not
let me go.

From this question came another, and then yet another. Each answer opened
the view to something I had not anticipated: a spiral, a cone, and finally
a space curve on which---as a corollary of the algebraic construction---all
prime numbers greater than $3$ appear as discrete points.

\medskip

This work follows the order in which these discoveries were made, not the
order of a retrospective deductive presentation. This is a deliberate choice.
In mathematics, one often distinguishes between the \emph{path of discovery}
and the \emph{path of exposition}---the former is inductive and tentative,
the latter deductive and orderly. Textbooks and dissertations almost always
choose the second path: they begin with axioms and definitions and proceed
to theorems as though someone already knew the answer before posing the
question.

This work begins with observation, develops conjectures from patterns, and
only then seeks the proof. At no point is the reader confronted with a result
whose origin remains unclear. Every concept appears when it is needed---not
before.

\medskip

Formally, the sequence under investigation is the arithmetic mean of the
first $n$ perfect squares:
\[
k(n) = \frac{1}{n}\sum_{j=1}^{n} j^2 = \frac{2n^2 + 3n + 1}{6}.
\]
That this formula carries this name and has this form was not known at the
outset. It was reconstructed from the table values---through differencing,
pattern recognition, and elementary algebra. The result of this reconstruction
is none other than the classical result found in every textbook on analysis.
The path to it, however, was a different one.

What carries this work beyond the mere derivation is the question:
\emph{What does this sequence mean geometrically?} The answer to this
question---an Archimedean spiral, a cone, a helix, a universal quantity
$Q(k) = \sqrt{48k+1}$ that unites all geometric coordinates in a single
expression---could not have been foreseen. It emerged step by step, as a
consequence of consistently pressing forward.

\medskip

Beyond its mathematical content, this work is also a testament to what
curiosity is capable of. Anyone willing to look closely can follow it---just
as it was created: step by step, number by number.

\bigskip

\noindent Dogan Balban\\
\noindent February 2026
